\chapter{Discussion}
\label{chap:discussion}

This chapter interprets the empirical findings of Chapter~\ref{chap:results} in broader context. Section~\ref{sec:disc_exchangeability} discusses the exchangeability assumption underlying conformal coverage and its relevance to the temporally ordered SPX options data. Section~\ref{sec:disc_practical} considers the practical interpretation of the resulting band surfaces. Section~\ref{sec:disc_limitations} discusses the limitations of the current study, and Section~\ref{sec:disc_future} outlines focused directions for future research. Section~\ref{sec:disc_summary} summarises the main points.

\section{The Exchangeability Assumption in Practice}
\label{sec:disc_exchangeability}

The standard finite-sample coverage guarantee of conformal prediction relies on exchangeability of the calibration and test nonconformity scores, meaning that their joint distribution is invariant to permutations of their indices \cite{shafer2008tutorial,angelopoulos2023gentle}. Independent and identically distributed observations provide a common sufficient condition, but exchangeability is the more general requirement. In the present application, option surfaces are ordered in time and the quotes within a surface are related through the same market state. Exact exchangeability is therefore not assumed. The rolling procedure is used as a practical adaptation mechanism, and its coverage is evaluated empirically in Chapter~\ref{chap:results}.

\subsection{Serial Dependence and the Rolling Window}
\label{subsec:disc_serial}

The rolling window contains the $L=5$ most recently processed surfaces. Because the available dates are irregular, it represents the latest five observations rather than a fixed calendar span. Older surfaces leave the window as new surfaces are processed. This restriction does not establish exchangeability or restore the formal finite-sample guarantee under temporal dependence. Its purpose is to make the conformal quantile responsive to recent errors, consistent with the motivation for giving recent observations greater relevance under distribution drift \cite{barber2023conformal}.

During the initial calibration period, the Score B global quantile is 0.016401 after the first surface and 0.016937 after surface 14. Across the 15 calibration steps, it ranges from 0.015899 to 0.018801. These values document how the threshold changes as surfaces enter and leave the window. They do not establish exchangeability or rule out distributional change; they only describe the behaviour of the selected rolling procedure before the walk-forward evaluation.

\subsection{Temporal Variation in Coverage}
\label{subsec:disc_regime}

Across the 31 evaluation surfaces, post-projection global coverage is below the 90\% target on ten surfaces, while Mondrian coverage is below 90\% on three. The three Mondrian shortfalls occur on surfaces 18, 28, and 42, corresponding to 17 April, 4 June, and 5 August 2026. For example, on surface 15 global coverage is 88.12\% and Mondrian coverage is 91.42\%, whereas surface 42 remains below the target at 88.84\% globally and 88.93\% with Mondrian assignment. The remaining shortfalls are therefore separated across the evaluation period rather than concentrated in one consecutive block.

At each walk-forward step, the current surface is evaluated before its scores are added to the window. The procedure can therefore respond only after a new error pattern has been observed; it cannot anticipate a change that is absent from the preceding surfaces. The available data do not include a separate analysis of market regimes or other covariates that would identify the cause of each coverage shortfall. Barber et al.\ \cite{barber2023conformal} analyse conformal coverage under departures from exchangeability. Accordingly, the reported coverage is interpreted as an empirical result, and the standard exchangeable-data guarantee is not claimed for this time-dependent evaluation.

\subsection{Contribution of Mondrian Binning}
\label{subsec:disc_mondrian_remedy}

Mondrian binning raises mean post-projection coverage from 90.80\% for the global construction to 92.85\%, an increase of 2.05 percentage points. On eight surfaces where global coverage is below 90\%, the Mondrian construction raises coverage to at least the target. It does not raise coverage on every surface, and surfaces 18, 28, and 42 remain below 90\%. The cross-surface standard deviation is 1.67 percentage points for global coverage and 1.97 percentage points for Mondrian coverage. The higher mean coverage therefore does not correspond to lower temporal variation.

At surface 14, the bin-specific Score B quantiles range from 0.011970 to 0.039470 around the global quantile of 0.016937. Each target point receives a quantile according to its $(\rho,z)$ region, while its final band width also depends on the local Vega weight and the no-arbitrage projection. Mondrian conformal prediction is motivated by calibration within predefined categories under exchangeability \cite{vovk2005algorithmic,shafer2008tutorial}. Here, the observed quantile pattern is evidence that the implementation adapts spatially; it is not a claim of exact conditional coverage.

\section{Practical Implications}
\label{sec:disc_practical}

The framework produces lower and upper surfaces around the smooth GNO estimate over the modelled domain. Calibration and empirical coverage are computed at observed quote locations, where market implied volatility is available for comparison. Interpolation and projection then convert the quote-level bands into surface objects. The output can therefore be examined across the domain without replacing the quote-level coverage calculation.

For Score B, the post-projection Mondrian bands have a mean IV width of 0.027180 and a mean half-width of 0.013590. The GNO mean absolute error on the same held-out target quotes is 0.006813. Band width and prediction error measure different quantities: the former is determined by the recent calibration scores, local Vega weighting, and projection, while the latter measures point-prediction accuracy. Wider regions indicate greater empirically calibrated prediction uncertainty around the GNO estimate. The bands summarise recent out-of-sample prediction errors under the selected score, rolling calibration, and projection; they are not confidence intervals for the GNO parameters or direct trading or hedging recommendations.

The comparison favours Score B for this experiment. After projection, its mean Mondrian coverage is 92.85\% with a mean IV width of 0.027180. The corresponding Score A coverage is 84.27\%, with a mean width of 0.067806, approximately 2.5 times the Score B width. Score A also has greater variation across evaluation surfaces: the standard deviation of post-projection Mondrian coverage is 6.24 percentage points, compared with 1.97 percentage points for Score B. Score B is therefore retained as the primary construction because it exceeds the nominal coverage level with narrower bands and less variation across surfaces in this evaluation.

Projection changes the Score B mean width by 0.000069, or 0.26\%, compared with an increase of 0.007008, or 9.88\%, for Score A. For both scores, the projection eliminates all violations of the enforced calendar and discrete butterfly conditions on the evaluation grids. No lower-to-upper crossings are observed at the evaluated quotes. These finite-grid results apply to the conditions enforced by the implementation and do not establish absence of arbitrage under every possible diagnostic.

\section{Limitations}
\label{sec:disc_limitations}

The empirical study uses one underlying asset, one pre-trained GNO checkpoint, and 56 stored SPX option surface snapshots. Ten snapshots with non-regular-session dates are excluded, leaving 46 retained surfaces. The first 15 surfaces, from 11 December 2025 to 17 February 2026, form the initial calibration period, and the remaining 31 surfaces, from 3 March to 11 August 2026, form the walk-forward evaluation. The saved date is the collection date recorded by the retrieval script rather than a separate exchange-provided quote timestamp. The observations are irregular and do not form a complete sample of trading days. The findings are therefore specific to this dataset, model, and period.

A development-only sensitivity check was conducted for Score B before projection. Surfaces 0 to 6 provide the initial history, while surfaces 7 to 14 form a common validation period of eight surfaces. One parameter is varied at a time: $L\in\{3,5,7\}$, a $3\times3$, $4\times4$, or $5\times5$ Mondrian partition, and $\lambda\in\{0.1,0.3,0.5\}$. The 31 final evaluation surfaces are not used. Across the seven configurations, mean coverage ranges from 89.79\% to 91.39\%, and mean IV width ranges from 0.02794 to 0.02868. The reference configuration $L=5$, $4\times4$, and $\lambda=0.3$ gives 91.11\% coverage and a mean width of 0.02831. Because the comparison uses one fixed context-target mask, only eight validation surfaces, and one-at-a-time parameter changes, it is a local sensitivity check rather than an exhaustive robustness or model-selection study.

Coverage is measured over the held-out target quotes within each evaluation surface and then averaged across surfaces. Quotes within a surface are cross-sectionally related, and regions containing more target quotes contribute more observations to that surface's coverage. The reported confidence intervals resample complete surfaces independently and do not model serial dependence between dates. Coverage and confidence intervals should therefore be interpreted as empirical summaries of the evaluation sample, not as exact finite-sample guarantees for the time series.

The projection enforces a fixed-$k$ calendar condition and a discrete butterfly condition based on neighbouring grid points on a finite $(\rho,z)$ grid. Their aggregate violation counts fall to zero after projection for both scores. A separate finite-difference evaluation of the GNO-style butterfly functional is retained as a supplementary diagnostic rather than an enforced constraint. Its post-projection counts remain nonzero: 7,432 for Score A and 2,731 for Score B across both bounds and all evaluation surfaces. The zero counts for the enforced conditions therefore do not imply that every arbitrage diagnostic is zero or that the bands are arbitrage-free over the continuous domain. Grid resolution and interpolation can also affect the reported counts.

The lower and upper surfaces are adjusted jointly, with their ordering included in the linear optimisation on the projection grid. No crossings are observed after interpolation at the evaluated quote locations for either score. This check does not examine every point of the continuous surface between the grid nodes and observed quotes.

\section{Future Research}
\label{sec:disc_future}

Future empirical work should evaluate the framework on a longer sequence of SPX surfaces, on additional index and single-name option datasets, and with more than one trained GNO checkpoint. Such experiments would test whether the relative performance of Score A and Score B, the observed band widths, and the effect of projection persist across market periods, underlyings, and model fits.

Temporal adaptation should be studied using market-state variables or predefined volatility regimes. Weighted or adaptive conformal methods could then be compared with the fixed rolling window when recent observations differ from older calibration data \cite{barber2023conformal}. The comparison should retain the walk-forward order so that information from a surface is unavailable until after that surface has been evaluated.

Parameter robustness requires a longer validation period, repeated context-target masks, and joint changes to the window length, Mondrian grid, and shrinkage parameter. Model selection and final evaluation should remain separate. Confidence intervals could also use a block or time-series resampling method that represents dependence between neighbouring dates rather than resampling surfaces independently.

\pagebreak
For the shape adjustment, the supplementary GNO-style butterfly functional could be incorporated directly into the projection objective or constraints and compared with the discrete butterfly condition currently enforced. Repeating the analysis on finer grids would help identify violations that occur between the current grid nodes. Comparisons with other uncertainty constructions should report coverage, band width, computational cost, enforced-condition counts, and the separate GNO-style butterfly diagnostic under the same evaluation masks.

\section{Summary}
\label{sec:disc_summary}

The walk-forward experiment contains 31 evaluation surfaces. For Score B, mean post-projection coverage increases from 90.80\% with global calibration to 92.85\% with Mondrian calibration, and the post-projection Mondrian bands have a mean IV width of 0.027180. The corresponding Score A construction has mean coverage of 84.27\% and a mean width of 0.067806. Its cross-surface coverage standard deviation is also higher, at 6.24 percentage points compared with 1.97 percentage points for Score B. Within this evaluation, Score B has higher coverage, narrower bands, and lower variation across surfaces.

For both scores, projection eliminates all violations of the enforced calendar and discrete butterfly conditions on the evaluation grids, and no lower-to-upper crossings are observed at the target quotes. The supplementary GNO-style butterfly diagnostic remains nonzero after projection, with aggregate counts of 7,432 for Score A and 2,731 for Score B. The conclusions are consequently limited to the implemented finite-grid conditions and the observed SPX sample. Temporal dependence, the use of one GNO checkpoint, and the small development-only sensitivity check also prevent the empirical results from being interpreted as universal coverage or arbitrage guarantees.
